% sections/08-testing-requirements.tex — Testing Requirements

\section{Testing Requirements}

% ── Overview ──
\subsection{Overview}

Comprehensive testing is mandatory. Your test suite must cover:

\begin{enumerate}
  \item \textbf{Unit tests} for each custom data structure
  \item \textbf{Integration tests} for module interactions
  \item \textbf{Edge case tests} for boundary conditions and error handling
  \item \textbf{Performance tests} for large inputs and stress scenarios
\end{enumerate}

% ── Test Organization ──
\subsection{Test Organization}

\begin{lstlisting}[style=tree]
tests/
|-- test_main.cpp .................. Test runner entry point
|-- unit/
|   |-- test_data_structures.cpp ... Data structure unit tests
|   |-- test_algorithms.cpp ........ Algorithm unit tests
|   +-- test_utils.cpp ............. Utility function tests
|-- integration/
|   |-- test_workflow.cpp .......... End-to-end workflow tests
|   +-- test_file_io.cpp .......... File I/O integration tests
|-- edge_cases/
|   |-- test_empty.cpp ............. Empty input tests
|   |-- test_single.cpp ........... Single element tests
|   +-- test_boundary.cpp ......... Boundary condition tests
|-- performance/
|   |-- test_stress.cpp ........... Stress tests with large inputs
|   +-- test_benchmark.cpp ........ Performance benchmarks
+-- fixtures/
    |-- test_helpers.hpp ........... Shared test utilities
    +-- mock_data.hpp .............. Mock data generators
\end{lstlisting}

% ── Test Coverage ──
\subsection{Test Coverage}

Every custom data structure must be tested across the following categories:

\begin{center}
\begin{tabular}{l l}
\hline
\rowcolor{tableheader} \textbf{Category} & \textbf{Test Cases} \\
\hline
Construction & Default constructor, parameterized constructor \\
Insertion & Beginning, middle, end, duplicate handling \\
Deletion & Beginning, middle, end, non-existent element \\
Search & Existing element, non-existent element \\
Edge Cases & Empty structure, single element, full capacity \\
Boundary & Maximum size, minimum values, overflow \\
\hline
\end{tabular}
\end{center}

% ── Test Format ──
\subsection{Test Format}

Each test should follow a clear structure with description, input, expected output, and verification:

\begin{lstlisting}[style=cpp]
void test_insert_at_beginning() {
    // Description: Insert element at position 0
    // Input: List [10, 20, 30], insert 5 at index 0
    // Expected: List [5, 10, 20, 30], size = 4

    LinkedList<int> list;
    list.insert(0, 10);
    list.insert(1, 20);
    list.insert(2, 30);

    // Precondition check
    assert(list.getSize() == 3);

    // Action
    list.insert(0, 5);

    // Postcondition verification
    assert(list.getSize() == 4);
    assert(list.get(0) == 5);
    assert(list.get(1) == 10);
    assert(list.get(2) == 20);
    assert(list.get(3) == 30);

    std::cout << "  test_insert_at_beginning PASSED\n";
}
\end{lstlisting}

% ── Memory Testing ──
\subsection{Memory Testing}

All pointer-based data structures must be tested for memory leaks:

\begin{lstlisting}[style=bash]
make memcheck
# or manually:
valgrind --leak-check=full --show-leak-kinds=all --track-origins=yes \
    ./build/bin/tests
\end{lstlisting}

Expected clean output:

\begin{lstlisting}[style=plaintext]
==12345== HEAP SUMMARY:
==12345==     in use at exit: 0 bytes in 0 blocks
==12345==   total heap usage: 150 allocs, 150 frees, 12,800 bytes allocated
==12345==
==12345== All heap blocks were freed -- no leaks are possible
\end{lstlisting}

% ── Test Frameworks ──
\subsection{Test Frameworks}

You may use any of the following testing approaches:

\begin{itemize}
  \item \textbf{Google Test} --- full-featured, industry standard
  \item \textbf{Catch2} --- header-only, modern \cpp
  \item \textbf{doctest} --- very fast compilation, minimal overhead
  \item \textbf{Custom assert-based testing} --- simple \code{assert()} with descriptive output
\end{itemize}

% ── Test Output ──
\subsection{Test Output}

Running \maketest{} should produce clear, structured output:

\begin{lstlisting}[style=plaintext]
[UNIT] Data Structures
  test_constructor .................. PASSED
  test_insert ....................... PASSED
  test_insert_at_beginning ......... PASSED
  test_insert_at_end ............... PASSED
  test_remove ...................... PASSED
  test_search ...................... PASSED
  test_size ........................ PASSED

[UNIT] Algorithms
  test_sort ........................ PASSED
  test_traverse .................... PASSED

[INTEGRATION] Workflow
  test_load_and_process ............ PASSED
  test_end_to_end .................. PASSED

[EDGE CASES] Boundary Conditions
  test_empty_structure ............. PASSED
  test_single_element .............. PASSED
  test_max_values .................. PASSED

[PERFORMANCE] Stress Tests
  test_large_input (10000 ops) ..... PASSED
  test_memory_usage ................ PASSED

SUMMARY: 16/16 tests passed
TIME:    0.042s
MEMORY:  No leaks detected
\end{lstlisting}

% ── Continuous Testing ──
\subsection{Continuous Testing}

Run your tests frequently throughout development:

\begin{itemize}
  \item After implementing each method
  \item Before committing changes
  \item After any refactoring
\end{itemize}
